% PILOTO INCREMENTAL V2 - revision conservadora y enriquecida.
% BASE: temas/tema1_piloto_incremental.tex.
% Se preservan: recorrido original, ejemplos, formulas, imagenes y ejercicios.
% Se han desarrollado diapositivas escuetas y figuras sin interpretacion.
% Dos diapositivas nuevas introducen los dominios con arcoseno/logaritmo
% y el calculo de niveles de una funcion racional; ambas refuerzan ejemplos.
% Cada marcador "Diapositiva original N" mantiene la trazabilidad.
% Archivo de prueba independiente: no altera tema1.tex, main.tex ni pilotos.
% PILOTO INCREMENTAL, basado directamente en temas/tema1.tex.
% Conserva la portada, las 21 diapositivas originales, sus figuras y ejercicios.
% Añade nueve diapositivas: índice, puentes, ejemplos y práctica graduada.
% Los marcadores "Diapositiva original N" garantizan trazabilidad.
% Archivo independiente: no se modifica el temario oficial ni otros pilotos.
% Compilar desde la raíz del repositorio.
\documentclass[main.tex]{subfiles}

\begin{document}
\graphicspath{{imagenes/tema1/}{../imagenes/tema1/}}

% Diapositiva original 1
\portadatema{1}{Funciones de varias variables}
% NUEVA: Índice general
\begin{frame}{Contenido del tema}
\begin{enumerate}
\item Funciones reales de varias variables.
\item Dominio, imagen y representación gráfica.
\item Determinación del dominio: ejemplos.
\item Operaciones con funciones.
\item Curvas y conjuntos de nivel.
\item Ejercicios de aplicación.
\end{enumerate}
\medskip
\small Los conceptos se ilustrarán también con aplicaciones de visualización científica e informática.
\end{frame}

% Diapositiva original 2
\begin{frame}{Funciones de varias variables}
La temperatura $T$ en un punto de la superficie de la Tierra a una hora determinada depende de la \textbf{longitud} y la \textbf{latitud} de dicho punto.
\medskip

Si fijamos el instante de observación y denominamos $x$ e $y$ a esas coordenadas, podemos expresar la temperatura como
\[
\boxed{T=f(x,y).}
\]
\medskip

Es una magnitud que depende de \textbf{dos variables independientes}. En general, cambiar la posición puede cambiar el valor de la temperatura.
\end{frame}
% Diapositiva original 3
\begin{frame}{Funciones de varias variables}
El volumen de una caja rectangular depende de sus tres dimensiones:
\[
\boxed{V(x,y,z)=x\cdot y\cdot z.}
\]
\begin{columns}[c]
\begin{column}{.54\textwidth}
Aquí las entradas son las dimensiones $x,y,z$ y la salida es el volumen.
\medskip

Para dimensiones físicas positivas: $x,y,z>0$.
\[
V(2,3,4)=24.
\]
\end{column}
\begin{column}{.44\textwidth}
\figura{s3_2.png}{.95}{.37}
\end{column}
\end{columns}
\end{frame}
% Diapositiva original 4
\begin{frame}{Funciones reales de $n$ variables}
\textbf{Definición.} Una función real de $n$ variables es una aplicación
\[
f:D\subseteq\mathbb R^n\longrightarrow\mathbb R
\]
que asigna a cada $n$-tupla $(x_1,\ldots,x_n)\in D$ un \textbf{único} número real:
\[
u=f(x_1,x_2,\ldots,x_n).
\]
\medskip

Para $n=2$, las entradas son pares $(x,y)$; para $n=3$, son ternas $(x,y,z)$. En ambos casos la salida es un número real.
\end{frame}
% Diapositiva original 5
\begin{frame}{De una variable a varias variables}
Los conceptos conocidos para funciones de una variable tienen su equivalente para funciones de $n$ variables.
\medskip

Las entradas pasan de ser números $x$ a vectores
$(x_1,\ldots,x_n)$, pero conservamos las ideas de
\textbf{dominio}, \textbf{imagen} y \textbf{gráfica}.
\medskip

Por ejemplo, para $f(x,y)=x^2+y^2$:
\[
D=\mathbb R^2,\qquad \operatorname{Im}(f)=[0,+\infty).
\]
El codominio puede ser $\mathbb R$, aunque la imagen sea menor.
\end{frame}
% Diapositiva original 6
\begin{frame}{Dominio e imagen de una función}
Para una función de dos variables $f:D\subseteq\mathbb R^2\to\mathbb R$:
\begin{itemize}
\item el \textbf{dominio} $D$ contiene los pares $(x,y)$ admitidos;
\item la \textbf{imagen} contiene los valores $f(x,y)$ que se obtienen;
\item el \textbf{codominio} se fija al definir la aplicación (aquí, $\mathbb R$).
\end{itemize}
\vspace{-1mm}
\figura{s6_0_v2.jpg}{.92}{.28}
\vspace{-2mm}
\small En el esquema, cada punto del dominio se asocia con \textbf{un valor real}. Distintos puntos pueden dar valores iguales.
\end{frame}
% Diapositiva original 7
\begin{frame}{Variables independientes y dependientes}
Las funciones de dos variables suelen escribirse como
\[
z=f(x,y),
\]
donde $x$ e $y$ son \textbf{variables independientes} y $z$ es la \textbf{variable dependiente}.
\medskip

Por ejemplo, si $f(x,y)=x+y^2$, entonces
\[
z=f(1,2)=1+2^2=5.
\]
\medskip

El par $(1,2)$ es la \textbf{entrada}; el número $5$ es su \textbf{imagen}. Al representar esta relación geométricamente, obtendremos el punto $(1,2,5)$.
\end{frame}
% NUEVA: Ejemplo de base antes de las gráficas
\begin{frame}{Ejemplo: evaluación y dominio}
Consideremos la función $f(x,y)=x^2+2y$.
\medskip

Evaluamos sustituyendo las dos variables:
\[
f(1,2)=1^2+2\cdot2=5,\qquad
f(-1,3)=(-1)^2+2\cdot3=7.
\]
\medskip

Su \textbf{dominio natural} es todo $\mathbb R^2$, porque no aparecen operaciones que impongan restricciones.
\medskip

Este ejemplo servirá como referencia al estudiar casos más complejos.
\end{frame}

% Diapositiva original 8
\begin{frame}{Representación gráfica de funciones de dos variables}
Para una función $f:D\subseteq\mathbb R^2\to\mathbb R$, su gráfica es
\[
G_f=\{(x,y,z)\in\mathbb R^3:\ (x,y)\in D,\ z=f(x,y)\}.
\]
\smallskip

Igual que las funciones de una variable se representan mediante curvas en el plano, las de dos variables se representan habitualmente mediante \textbf{superficies en el espacio}.
\vspace{-2mm}
\begin{columns}[c]
\begin{column}{.49\textwidth}
\figura{s8_2.png}{.97}{.40}
\end{column}
\begin{column}{.49\textwidth}
\small Cada punto $(x,y)$ del dominio se corresponde con el punto de altura $z=f(x,y)$.
\medskip

\textbf{No debe confundirse} el dominio, que está en el plano $XY$, con la gráfica, que está en $\mathbb R^3$.
\end{column}
\end{columns}
\end{frame}
% Diapositiva original 9
\begin{frame}{Ejemplo: superficie en el espacio}
Consideremos $z=f(x,y)=x+y^2$, cuyo dominio es todo $\mathbb R^2$.
\medskip

Fijado $y$, la función varía linealmente con $x$; fijado $x$, su dependencia respecto de $y$ es cuadrática.
\vspace{-2mm}
\figura{s9_2.png}{.82}{.45}
\vspace{-3mm}
\small Es una superficie que combina comportamientos distintos según la dirección considerada.
\end{frame}
% NUEVA: Contexto para los ejemplos originales de dominio
\begin{frame}{Determinación del dominio natural}
El ejemplo polinómico anterior tiene dominio $\mathbb R^2$.
En otras expresiones, las operaciones introducen restricciones.
\medskip
\begin{itemize}
\item En un \textbf{cociente}, el denominador no puede ser cero.
\item En una \textbf{raíz cuadrada}, el radicando debe ser no negativo.
\item En un \textbf{logaritmo}, el argumento debe ser positivo.
\end{itemize}
\medskip

Si aparecen varias condiciones, el dominio es su
\textbf{intersección}. Veamos tres ejemplos.
\end{frame}

% Diapositiva original 10
\begin{frame}{Ejemplo: dominio de una función racional}
En la función
\[
f(x,y)=\frac{3x-5y}{y-x^2}
\]
la única restricción proviene del denominador:
\[
y-x^2\ne0 \quad\Longleftrightarrow\quad y\ne x^2.
\]
\[
D=\{(x,y)\in\mathbb R^2:y\ne x^2\}.
\]
\vspace{-2mm}
\figura{s10_4.png}{.66}{.29}
\vspace{-2mm}
\small La parábola $y=x^2$ queda excluida del dominio. Al aproximarnos a ella pueden aparecer valores muy grandes de la función.
\end{frame}
% Diapositiva original 11
\begin{frame}{Ejemplo: dominio de una función radical}
En la función
\[
f(x,y)=\sqrt{9-x^2-y^2}
\]
el radicando debe ser no negativo:
\[
9-x^2-y^2\ge0\ \Longleftrightarrow\ x^2+y^2\le9.
\]
\[
D=\{(x,y)\in\mathbb R^2:x^2+y^2\le9\}.
\]
\vspace{-2mm}
\figura{s11_5.png}{.66}{.28}
\vspace{-2mm}
\small El dominio es el disco cerrado de radio $3$; la gráfica es la semiesfera superior.
\end{frame}
% Diapositiva original 12
\begin{frame}{Ejemplo: dominio de una función de tres variables}
Consideremos
\[
f(x,y,z)=\ln(z-y)+xy\sin z.
\]
\medskip

El término $xy\sin z$ está definido para todo $(x,y,z)\in\mathbb R^3$. El logaritmo exige, en cambio,
\[
z-y>0\quad\Longleftrightarrow\quad z>y.
\]
Por tanto,
\[
D=\{(x,y,z)\in\mathbb R^3:z>y\}.
\]
\medskip

Geométricamente, es el \textbf{semiespacio abierto} situado del lado $z>y$ del plano $z=y$. El propio plano queda excluido.
\end{frame}
% NUEVA: Dominio con restricciones combinadas
\begin{frame}{Ejemplo adicional: dos restricciones simultáneas}
Consideremos
\[
f(x,y)=\frac{\sqrt{4-x^2-y^2}}{y-1}.
\]
Debe cumplirse simultáneamente:
\[
4-x^2-y^2\geq0,\qquad y-1\ne0.
\]
Por tanto,
\[
\boxed{D=\{(x,y)\in\mathbb R^2:
x^2+y^2\leq4,\ y\ne1\}.}
\]
\medskip

Es el disco cerrado de radio $2$, del que se excluyen los puntos de la recta $y=1$ que pertenecen al disco.
\end{frame}


% NUEVA V2: Ejemplo puente para el ejercicio avanzado del final
\begin{frame}{Ejemplo adicional: arcoseno y logaritmo}
Para determinar el dominio de
\[
g(x,y,z)=\arcsin(xy)+\ln(x+z)
\]
debemos imponer dos condiciones:
\[
-1\le xy\le1
\qquad\text{y}\qquad x+z>0.
\]
\medskip

Por tanto,
\[
D=\{(x,y,z)\in\mathbb R^3:
-1\le xy\le1,\ x+z>0\}.
\]
\medskip

Este procedimiento se aplica también a expresiones más complejas: \textbf{cada operación aporta sus propias restricciones}.
\end{frame}

% Diapositiva original 13
\begin{frame}{Operaciones con funciones de varias variables}
Las funciones de varias variables pueden operarse de igual forma que las de una variable, ya sea con suma, producto, cociente o composición de funciones.
\medskip

Para $f(x,y)=x+y$ y $g(x,y)=x^2+y^2$:
\[
(f+g)(x,y)=x+y+x^2+y^2,
\qquad
\left(\frac fg\right)(x,y)=\frac{x+y}{x^2+y^2}.
\]
La suma está definida en $\mathbb R^2$;
en el cociente hay que excluir $(0,0)$, donde $g=0$.
\end{frame}
% NUEVA: Puente conceptual hacia las curvas de nivel
\begin{frame}{De la superficie a la representación plana}
Las gráficas de funciones de dos variables se representan generalmente mediante superficies tridimensionales.
\medskip

En cartografía, representación de datos y visualización científica
interesa también describirlas mediante dibujos en el plano.
\medskip

Una técnica consiste en representar los puntos del dominio donde la función toma \textbf{el mismo valor constante}.
\medskip

Esta idea conduce a las \textbf{curvas de nivel}.
\end{frame}

% Diapositiva original 14
\begin{frame}{Curvas de nivel: definición e interpretación}
Las \textbf{curvas de nivel} de $f:D\subseteq\mathbb R^2\to\mathbb R$ son los conjuntos que satisfacen
\[
C_k=\{(x,y)\in D:f(x,y)=k\},
\]
donde $k$ es una constante.
\vspace{-1mm}
\begin{columns}[c]
\begin{column}{.48\textwidth}
\figura{s14_3.png}{.98}{.39}
\end{column}
\begin{column}{.48\textwidth}
\figura{s14_5.png}{.98}{.39}
\end{column}
\end{columns}
\vspace{-3mm}
\small Izquierda: corte horizontal y proyección de un nivel sobre el plano. Derecha: una superficie con varios máximos y mínimos también admite curvas de nivel.
\end{frame}
% NUEVA: Primer ejemplo de curvas de nivel sencillo
\begin{frame}{Ejemplo adicional: niveles de una función afín}
Sea $f(x,y)=x+y$. Para cada nivel $k$:
\[
f(x,y)=k\quad\Longleftrightarrow\quad y=k-x.
\]
Las curvas de nivel son \textbf{rectas paralelas}.
\[
\begin{array}{c|c}
k & \text{curva de nivel}\\\hline
-1 & y=-1-x\\
0 & y=-x\\
1 & y=1-x
\end{array}
\]
\medskip

El siguiente ejemplo original permite estudiar una geometría diferente: las circunferencias.
\end{frame}

% Diapositiva original 15
\begin{frame}{Ejemplo: curvas de nivel de una semiesfera}
\textbf{Ejemplo original.} Obtener y clasificar las curvas de nivel de
\[
z=f(x,y)=\sqrt{100-x^2-y^2}.
\]
\medskip

El dominio es el disco $x^2+y^2\le100$ y la imagen es $[0,10]$.
\medskip

Fijaremos una altura $z=c$ y calcularemos todos los puntos $(x,y)$ que la alcanzan.
\medskip

Así relacionaremos el \textbf{valor de la función} con la \textbf{geometría} de sus curvas de nivel.
\end{frame}
% Diapositiva original 16
\begin{frame}{Resolución: niveles de la semiesfera}
Para $f(x,y)=\sqrt{100-x^2-y^2}$, fijamos la altura $c$:
\begin{align*}
c&=\sqrt{100-x^2-y^2},\\
c^2&=100-x^2-y^2,\\
x^2+y^2&=100-c^2.
\end{align*}
\medskip

\textbf{Interpretación según el nivel:}
\begin{itemize}
\item $0\le c<10$: circunferencia de radio $\sqrt{100-c^2}$;
\item $c=10$: únicamente el punto $(0,0)$;
\item $c<0$ o $c>10$: conjunto vacío.
\end{itemize}
\smallskip
\small La condición $c\ge0$ es esencial: no debe perderse al elevar al cuadrado.
\end{frame}
% Diapositiva original 17
\begin{frame}{Representación gráfica de los niveles}
En el ejemplo $z=\sqrt{100-x^2-y^2}$, cada plano horizontal $z=c$ corta la semiesfera en una circunferencia (si $0\le c<10$).
\begin{columns}[c]
\begin{column}{.43\textwidth}
\[
x^2+y^2=100-c^2.
\]
\small
Al aumentar $c$, el radio $\sqrt{100-c^2}$ disminuye.
\medskip

En $c=10$ solo queda la cima.
\end{column}
\begin{column}{.55\textwidth}
\figura{s17_3.png}{1}{.45}
\end{column}
\end{columns}
\end{frame}
% Diapositiva original 18
\begin{frame}{Relación entre superficie y curvas de nivel}
Una curva de nivel es el resultado de cortar una superficie mediante un \textbf{plano horizontal} y proyectar el corte sobre el plano $XY$.
\begin{columns}[c]
\begin{column}{.65\textwidth}
\figura{s18_0.jpg}{.99}{.59}
\end{column}
\begin{column}{.33\textwidth}
\small
En la figura, cada línea horizontal sobre la superficie corresponde a un valor de $k$.
\medskip

Su proyección permite representar las mismas alturas sin dibujar la superficie 3D.
\end{column}
\end{columns}
\end{frame}
% Diapositiva original 19
\begin{frame}{Curvas de nivel de una función no radial}
Una superficie no tiene por qué presentar simetría circular. Consideremos
\[
f(x,y)=-xy\,e^{-x^2-y^2}.
\]
\vspace{-1mm}
\figura{s19_1.jpg}{.96}{.48}
\vspace{-3mm}
\small Como $e^{-x^2-y^2}>0$, el signo de $f$ viene determinado por $-xy$: es positivo si $xy<0$, negativo si $xy>0$, y nulo sobre los ejes.
\end{frame}
% Diapositiva original 20
\begin{frame}{Curvas de nivel de una función racional}
Para la función
\[
f(x,y)=\frac{-3y}{x^2+y^2+1}
\]
el denominador siempre es positivo, de modo que el signo de $f$ es el contrario del de $y$.
\vspace{-2mm}
\figura{s20_1.jpg}{.94}{.47}
\vspace{-3mm}
\small Las curvas de nivel reflejan una región de valores positivos ($y<0$) y otra de valores negativos ($y>0$). El nivel $0$ es el eje $x$.
\end{frame}

% NUEVA V2: Derivación algebraica de las curvas de la imagen original
\begin{frame}{Cálculo de los niveles de la función racional}
Para
\[
f(x,y)=\frac{-3y}{x^2+y^2+1},
\]
fijamos un nivel $k\ne0$:
\[
k(x^2+y^2+1)=-3y.
\]
Completando cuadrados:
\[
\boxed{x^2+\left(y+\frac{3}{2k}\right)^2
=\frac{9}{4k^2}-1.}
\]
\smallskip
\small
Son circunferencias si $0<|k|<\frac32$; un punto si $|k|=\frac32$;
y no hay puntos si $|k|>\frac32$.
Para $k=0$, el conjunto de nivel es la recta $y=0$.
\end{frame}

% Diapositiva original 21
\begin{frame}{Aplicación: curvas de nivel en cartografía}
Las curvas de nivel se emplean para representar en un plano los puntos que tienen la \textbf{misma altitud}.
\begin{columns}[c]
\begin{column}{.57\textwidth}
\figura{s21_0.jpg}{.99}{.66}
\end{column}
\begin{column}{.41\textwidth}
\small
Cada curva lleva asociada una cota (por ejemplo, 4500, 5000 o 5500).
\medskip

Cuando las curvas aparecen \textbf{muy próximas}, la altitud cambia rápidamente con la distancia horizontal: la pendiente es mayor.
\medskip

Este principio es la base de los mapas topográficos.
\end{column}
\end{columns}
\end{frame}
% NUEVA: Informática sin cambiar el hilo matemático
\begin{frame}{Aplicación informática: funciones de coste}
En optimización y aprendizaje automático, una función de coste depende de los parámetros de un modelo.
\medskip

Por ejemplo:
\[
L(w_1,w_2)=(w_1-2)^2+2(w_2+1)^2.
\]
Sus curvas de nivel verifican
\[
(w_1-2)^2+2(w_2+1)^2=k.
\]
\medskip

Para $k>0$ son \textbf{elipses} centradas en $(2,-1)$.
Estas curvas permiten interpretar cómo varía el coste al modificar los parámetros.
\end{frame}

% NUEVA: Práctica complementaria básica
\begin{frame}{Ejercicios complementarios: nivel básico}
\begin{enumerate}
\item Si $f(x,y)=2x+y$, calcular $f(1,2)$ y $f(-1,3)$.
\medskip
\item Determinar el dominio natural de
\[
g(x,y)=\sqrt{9-x^2-y^2}.
\]
\item Describir las curvas de nivel $k=0,1$ de
\[
h(x,y)=x+y.
\]
\end{enumerate}
\medskip
\small Estos ejercicios trabajan directamente los conceptos y procedimientos elementales.
\end{frame}

% NUEVA: Práctica complementaria intermedia
\begin{frame}{Ejercicios complementarios: nivel intermedio}
\begin{enumerate}
\item Determinar y dibujar el dominio de $f(x,y)=\ln(y-x^2)$.
\medskip
\item Encontrar el dominio de
\[
g(x,y)=\frac{\sqrt{4-x^2-y^2}}{y-1}.
\]
\item Describir geométricamente los niveles $k=0,1,4$ de
\[
h(x,y)=x^2+y^2.
\]
\end{enumerate}
\end{frame}

% Diapositiva original 22
\begin{frame}{Ejercicios originales: consolidación y ampliación}
\small\textbf{Ejercicios}
\medskip

Encontrar el dominio de la función
\[f(x,y,z)=\arcsin(xy)\,e^{3z}+e^{(y+z)\ln(x+z)}.\]
Describir los conjuntos de nivel para los valores de $k$ indicados:
\begin{enumerate}[a)]
\item $f(x,y)=x^2+y^2$,\quad $k=0,1,2,3$.
\item $f(x,y,z)=x^2+y^2$,\quad $k=0,1,2$.
\item $f(x,y)=x^2-y^2$,\quad $k=-2,-1,0,1,2$.
\end{enumerate}
En a) y c) se trata de curvas de nivel; en b), de superficies de nivel.
\end{frame}
\end{document}
